%%%%%%%%%%%%%%%%%%%%%%%%%%%%%%%%%%%%%%%%%%%%%%%%%%%%%%%%%%%%%%%%%%%%%%%%%%%%%%%%
% 08_measurement.tex: Precision measurement of the Higgs boson mass:
%%%%%%%%%%%%%%%%%%%%%%%%%%%%%%%%%%%%%%%%%%%%%%%%%%%%%%%%%%%%%%%%%%%%%%%%%%%%%%%%
\chapter{Higgs boson mass measurement}
\label{chap:measurement}
%%%%%%%%%%%%%%%%%%%%%%%%%%%%%%%%%%%%%%%%%%%%%%%%%%%%%%%%%%%%%%%%%%%%%%%%%%%%%%%%

%%%%%%%%%%%%%%%%%%%%%%%%%%%%%%%%%%%%%%%%%%%%%%%%%%%%%%%%%%%%%%%%%%%%%%%%%%%%%%%%

The expected mass of the Higgs boson in the diphoton decay channel is measured to be 
\begin{equation*}
    m_{\mathrm{H}} = 125.574 \pm 0.088 (\mathrm{stat.}) \pm 0.121 (\mathrm{syst.}) \mathrm{\ GeV}.
\end{equation*}

The Higgs mass likelihood profile for the combined measurement is shown in Figure~\ref{fig:mass_scans}. This measurement is preliminary; the CMS Collaboration will be publishing a final measurement in the coming months. The precision of the Higgs mass measurement in the diphoton channel is limited by the systematic uncertainty. Previously, the mass in the diphoton channel was measured by the CMS Collaboration using the 2016 dataset to be $m_{\mathrm{H}} = 125.78 \pm 0.18 (\mathrm{stat.}) \pm 0.18 (\mathrm{syst.})$ GeV~\cite{cms:2020:higgs_mass}, and the ATLAS Collaboration measured the mass in the diphoton channel using the full Run 2 dataset to be $m_{\mathrm{H}} = 125.17 \pm 0.11 (\mathrm{stat.}) \pm 0.09 (\mathrm{syst.})$ GeV~\cite{atlas:2023:higgs_mass}. The measurement presented in this thesis is the most precise measurement of the Higgs boson mass in the diphoton channel by the CMS Collaboration to date.

\begin{figure}[htbp!]
    \centering
    \includegraphics[width=0.75\linewidth]{08-figures/higgs_mass.png}
    \caption{The negative log-likelihood scans for the best-fit measurement of the Higgs mass are shown. The dotted line shows the ``stat. only" scan where the systematic uncertainties have been frozen in the fit, the solid line shows the full scan.}
    \label{fig:mass_scans}
\end{figure}

